\subsection{Durchgeführte Projekte}\label{subsec:DurchgeführteProjekte}

\par Wie in \fullref{subsec:Ausgangssituation} beschrieben, befindet sich die \gls{stz} aktuell in einer frühen Phase der Unternehmensentwicklung. Nichtsdestotrotz wurde mit der Generalsanierung der Liegenschaft im Marterlweg 18 (2024-2025) in Pollanten bereits ein erstes sehr umfangreiches Projekt erfolgreich abgeschlossen. Die Sanierung umfasste die Entkernung der bestehenden Immobilie, die Umwandlung von einem Einfamilienhaus in ein Mehrfamilienhaus für drei Parteien. Dies umfasste eine Grundrissänderung, den Anbau eines Treppenhauses, die umfassende Modernisierung der Haustechnik, die Erneuerung der Fassade, einen Neubau einer Doppelgarage und die Neugestaltung der Außenanlagen. Das Projekt wurde innerhalb des geplanten Zeit- und Budgetrahmens abgeschlossen und hat zu einer deutlichen Wertsteigerung der Immobilie beigetragen.

\par Auch die weiteren Liegenschaften der \gls{stz} wurden vor der Einbringung jeweils unter Federführung des Komplementärs und der Kommanditisten umfassend saniert, modernisiert, instandgesetzt oder wie im Falle der Hopfengasse 25A in Sollngriesbach (1994) und der Liegeneschaft im Marterlweg 18 in Pollanten (1988) von Grundauf neu errichtet. Umfassende Sanierungsarbeiten wurden in den Liegenschaften  Hierbei wurden insbesondere die Wohn- und Geschäftshaus in Pollanten (Tannenstraße 5) sowie die Doppelhaushälfte in Sollngriesbach (Hopfengasse 25A) einer umfangreichen Sanierung unterzogen, um den Wohnkomfort zu erhöhen, die Energieeffizienz zu verbessern und den Wert der Immobilien zu steigern. Auch in der Liegenschaft in München (Silcherstraße 8) wurden umfassende Sanierungsarbeiten durchgeführt, um die Wohnheinheit auf den neuesten Stand der Technik zu bringen und den Wohnkomfort für die Mieter zu erhöhen. Darüberhinaus wurde 2017 in der Liegenschaft in München (Silcherstraße 8) von der Hausgemeinschaft eine umfassende energetische Sanierung durchgeführt, um die Energieeffizienz zu verbessern und die Betriebskosten zu senken. Hierbei wurden insbesondere die Fassadendämmung, das Dach, die Fenster der Immobilie modernisiert, um den Energieverbrauch zu reduzieren und die Umweltbelastung zu minimieren. Derzeit findet eine Sanierung der zur Wohnanlage gehörenden Tiefgarage statt. Das Engagement der \gls{stz} in Form des Komplementärs im Verwaltungsbeirat der Wohnanlage zeigt das hohe Engagement der \gls{stz} für  die Sanierung und Modernisierung ihrer Liegenschaften. Dies spiegelt das langfristige Ziel wider, ein attraktives Immobilienportfolio aufzubauen und zu verwalten, um langfristige Wertsteigerungen zu erzielen. Durch die kontinuierliche Investition in die Instandhaltung und Modernisierung der bestehenden Immobilien trägt die \gls{stz} dazu bei, den Wert ihrer Liegenschaften zu erhalten und zu steigern, um langfristig eine nachhaltige Rendite für die Gesellschafter zu erzielen.

\par Insofern wurden bisher für alle Liegenschaften bereits Immobilienprojekte im klassischen Sinne durchgeführt. In den kommenden Jahren plant die \gls{stz} jedoch, weitere Projekte durchzuführen, insbesondere im Bereich des Immobilienerwerbs und der Entwicklung von Immobilien, um das Portfolio zu erweitern und langfristige Wertsteigerungen zu erzielen. Hierbei werden insbesondere Projekte im Süddeutschen Raum angestrebt, um von der positiven Entwicklung des Immobilienmarktes in der Region zu profitieren.